% ============================================================
%  Chapter 1 — Introduction & Background
%  Thesis: Detection of the pathogenicity in the clinical
%          microbial Nanopore sequencing
%
%  HOW TO USE:
%  - Paste this file into your Overleaf project
%  - Add % ============================================================
%  Chapter 1 — Introduction & Background
%  Thesis: Detection of the pathogenicity in the clinical
%          microbial Nanopore sequencing
%
%  HOW TO USE:
%  - Paste this file into your Overleaf project
%  - Add % ============================================================
%  Chapter 1 — Introduction & Background
%  Thesis: Detection of the pathogenicity in the clinical
%          microbial Nanopore sequencing
%
%  HOW TO USE:
%  - Paste this file into your Overleaf project
%  - Add % ============================================================
%  Chapter 1 — Introduction & Background
%  Thesis: Detection of the pathogenicity in the clinical
%          microbial Nanopore sequencing
%
%  HOW TO USE:
%  - Paste this file into your Overleaf project
%  - Add \input{introduction} in your main.tex
%  - Replace every TODO comment with the real content
%  - Add your bibliography entries (BibTeX keys are shown next to \cite{})
% ============================================================

\chapter{Introduction}

\section{Clinical Microbiology and Infectious Disease Diagnosis}

Infectious diseases remain a leading cause of morbidity and mortality worldwide,
accounting for a significant proportion of hospital admissions and intensive care
unit stays~\cite{TODO_WHO_infectious_disease}. Rapid and accurate identification
of the causative pathogen is essential for guiding antibiotic therapy, preventing
the spread of infection, and improving patient outcomes.

Traditional diagnostic workflows in clinical microbiology rely on culture-based
methods, which, despite their established reliability, suffer from several
limitations. Bacterial culture typically requires 24 to 72 hours to yield a
result, and many clinically relevant organisms are fastidious or entirely
unculturable under standard laboratory conditions~\cite{TODO_culture_limitations}.
During this waiting period, clinicians are often forced to prescribe broad-spectrum
antibiotics empirically, a practice that contributes to the global problem of
antimicrobial resistance (AMR)~\cite{TODO_AMR_global}.

\section{Next-Generation Sequencing in Clinical Practice}

Next-generation sequencing (NGS) technologies have transformed microbiology by
enabling the direct sequencing of genetic material from clinical samples, bypassing
the need for culture. Metagenomic sequencing in particular allows the simultaneous
detection of bacteria, viruses, fungi, and parasites from a single sequencing run,
providing a comprehensive view of the microbial content of a sample~\cite{TODO_metagenomics_review}.

Short-read sequencing platforms, such as Illumina, dominate the current NGS
landscape. However, the short read lengths produced by these platforms
(typically 150--300 base pairs) pose challenges for resolving complex genomic
regions, mobile genetic elements, and structural variants that are often associated
with virulence and antibiotic resistance~\cite{TODO_shortread_limitations}.

\section{Oxford Nanopore Sequencing}

Oxford Nanopore Technology (ONT) is a third-generation sequencing platform that
produces long reads — typically ranging from several kilobases to over one megabase
in length — by measuring changes in ionic current as a DNA or RNA strand passes
through a nanoscale protein pore~\cite{TODO_ONT_original}. This approach confers
several advantages that are particularly relevant in a clinical context:

\begin{itemize}
    \item \textbf{Speed.} Sequencing results can be obtained within hours of sample
          collection, enabling real-time clinical decision-making~\cite{TODO_ONT_speed}.
    \item \textbf{Long reads.} Read lengths of tens of kilobases allow complete
          assembly of plasmids and mobile genetic elements carrying resistance genes,
          which are often fragmented or misassembled by short-read technologies.
    \item \textbf{Portability and low cost.} The MinION device, weighing less than
          100 grams, can be operated at the point of care, in field settings, or in
          resource-limited environments where large sequencing infrastructure is
          unavailable~\cite{TODO_MinION_pointofcare}.
\end{itemize}

Despite these advantages, Nanopore sequencing is characterised by a higher per-base
error rate (approximately 5--10\% in recent chemistry generations) compared to
Illumina~\cite{TODO_nanopore_accuracy}. This must be accounted for in downstream
bioinformatics analyses, for example by using alignment parameters tuned for
long reads with elevated mismatch rates.

\section{Pathogenicity and Virulence}

Pathogenicity is defined as the ability of a microorganism to cause disease in a
susceptible host. It is not a binary property but rather a complex phenotype
determined by the combined action of multiple genes distributed across the
genome~\cite{TODO_pathogenicity_definition}.

Virulence factors are the molecular mechanisms through which a pathogen causes
damage to the host. These include, but are not limited to:

\begin{itemize}
    \item \textbf{Adhesins and pili} — enable the bacterium to attach to host tissues.
    \item \textbf{Toxins (exotoxins and endotoxins)} — disrupt host cell function or
          trigger damaging immune responses.
    \item \textbf{Secretion systems (Type I–VI)} — inject effector proteins directly
          into host cells to subvert immune defences.
    \item \textbf{Iron acquisition systems} — allow the pathogen to scavenge iron from
          the host, which is essential for growth.
    \item \textbf{Capsules and immune evasion factors} — protect the bacterium from
          phagocytosis and complement-mediated killing.
    \item \textbf{Biofilm formation} — enables persistent colonisation on tissues and
          implanted medical devices.
\end{itemize}

Because virulence is encoded across many, often distant genomic loci, and because
individual genes may be horizontally transferred between species on mobile genetic
elements, detecting pathogenicity requires a genome-wide analysis rather than
the identification of a single marker gene~\cite{TODO_HGT_virulence}.

\section{Virulence Factor and Antibiotic Resistance Databases}

Two curated public databases underpin the computational detection of virulence
and resistance genes:

\textbf{VFDB (Virulence Factor Database)}~\cite{TODO_VFDB_citation} is a manually curated
repository of virulence-associated sequences from experimentally verified pathogenic
organisms. It is organised into functional categories including adherence, motility,
toxin production, secretion systems, iron acquisition, immune evasion, biofilm, and
invasion. In this work, VFDB was used as the primary reference for virulence gene
detection.
% TODO: Add exact download date and sequence count here, e.g.:
% The version used in this work was downloaded on [DATE] and contained
% [N] protein sequences across [M] functional categories.

\textbf{CARD (Comprehensive Antibiotic Resistance Database)}~\cite{TODO_CARD_citation}
provides a curated collection of antibiotic resistance genes (ARGs) and their
associated mechanisms. CARD is used in this work to assess the antimicrobial
resistance potential of the sequenced microbiome, which is a key component of
clinical pathogenicity assessment.

\section{Computational Detection of Pathogenicity from Nanopore Data}

The bioinformatics workflow for detecting virulence genes from Nanopore sequencing
data involves several steps. First, gene-coding regions are predicted on the raw
long reads using tools such as Prodigal~\cite{TODO_prodigal} or its Python
bindings, Pyrodigal. Predicted coding sequences are then translated into protein
sequences and aligned against the VFDB protein database using DIAMOND
\textit{blastx}~\cite{TODO_diamond}, which natively handles the translation step
from nucleotide queries to protein references.

The resulting alignments are evaluated according to several quality metrics:
percent identity, alignment length, query coverage, and bit score. Thresholds on
these metrics are selected based on the observed score distributions (see
Chapter~\ref{ch:methods}) to distinguish biologically meaningful hits from
spurious low-quality alignments.

Virulence hits are then grouped into functional categories and quantified using
TPM (Transcripts Per Million) normalisation, which allows abundance comparisons
across samples with different total sequencing depths~\cite{TODO_TPM}.

\section{Aims of This Work}

The primary aim of this thesis is to develop a computational pipeline that:

\begin{enumerate}
    \item Accepts Nanopore long-read sequencing data (FASTA format) as input.
    \item Detects and quantifies virulence-associated genes by alignment against VFDB.
    \item Detects and quantifies antibiotic resistance genes by alignment against CARD.
    \item Classifies the detected genes into functional pathogenicity categories.
    \item Assigns a clinically interpretable risk level to each sample.
    \item Enables comparative analysis across multiple samples to assess changes in
          virulence content over time or across conditions.
\end{enumerate}

A secondary aim is to explore the feasibility of using large language models (LLMs)
and protein language models as an alternative, alignment-free approach to virulence
classification, and to compare its performance against the alignment-based pipeline.

The remainder of this thesis is structured as follows.
Chapter~\ref{ch:methods} describes the materials and methods.
% TODO: add remaining chapter descriptions once you have the structure in your main.tex
%  - Replace every TODO comment with the real content
%  - Add your bibliography entries (BibTeX keys are shown next to \cite{})
% ============================================================

\chapter{Introduction}

\section{Clinical Microbiology and Infectious Disease Diagnosis}

Infectious diseases remain a leading cause of morbidity and mortality worldwide,
accounting for a significant proportion of hospital admissions and intensive care
unit stays~\cite{TODO_WHO_infectious_disease}. Rapid and accurate identification
of the causative pathogen is essential for guiding antibiotic therapy, preventing
the spread of infection, and improving patient outcomes.

Traditional diagnostic workflows in clinical microbiology rely on culture-based
methods, which, despite their established reliability, suffer from several
limitations. Bacterial culture typically requires 24 to 72 hours to yield a
result, and many clinically relevant organisms are fastidious or entirely
unculturable under standard laboratory conditions~\cite{TODO_culture_limitations}.
During this waiting period, clinicians are often forced to prescribe broad-spectrum
antibiotics empirically, a practice that contributes to the global problem of
antimicrobial resistance (AMR)~\cite{TODO_AMR_global}.

\section{Next-Generation Sequencing in Clinical Practice}

Next-generation sequencing (NGS) technologies have transformed microbiology by
enabling the direct sequencing of genetic material from clinical samples, bypassing
the need for culture. Metagenomic sequencing in particular allows the simultaneous
detection of bacteria, viruses, fungi, and parasites from a single sequencing run,
providing a comprehensive view of the microbial content of a sample~\cite{TODO_metagenomics_review}.

Short-read sequencing platforms, such as Illumina, dominate the current NGS
landscape. However, the short read lengths produced by these platforms
(typically 150--300 base pairs) pose challenges for resolving complex genomic
regions, mobile genetic elements, and structural variants that are often associated
with virulence and antibiotic resistance~\cite{TODO_shortread_limitations}.

\section{Oxford Nanopore Sequencing}

Oxford Nanopore Technology (ONT) is a third-generation sequencing platform that
produces long reads — typically ranging from several kilobases to over one megabase
in length — by measuring changes in ionic current as a DNA or RNA strand passes
through a nanoscale protein pore~\cite{TODO_ONT_original}. This approach confers
several advantages that are particularly relevant in a clinical context:

\begin{itemize}
    \item \textbf{Speed.} Sequencing results can be obtained within hours of sample
          collection, enabling real-time clinical decision-making~\cite{TODO_ONT_speed}.
    \item \textbf{Long reads.} Read lengths of tens of kilobases allow complete
          assembly of plasmids and mobile genetic elements carrying resistance genes,
          which are often fragmented or misassembled by short-read technologies.
    \item \textbf{Portability and low cost.} The MinION device, weighing less than
          100 grams, can be operated at the point of care, in field settings, or in
          resource-limited environments where large sequencing infrastructure is
          unavailable~\cite{TODO_MinION_pointofcare}.
\end{itemize}

Despite these advantages, Nanopore sequencing is characterised by a higher per-base
error rate (approximately 5--10\% in recent chemistry generations) compared to
Illumina~\cite{TODO_nanopore_accuracy}. This must be accounted for in downstream
bioinformatics analyses, for example by using alignment parameters tuned for
long reads with elevated mismatch rates.

\section{Pathogenicity and Virulence}

Pathogenicity is defined as the ability of a microorganism to cause disease in a
susceptible host. It is not a binary property but rather a complex phenotype
determined by the combined action of multiple genes distributed across the
genome~\cite{TODO_pathogenicity_definition}.

Virulence factors are the molecular mechanisms through which a pathogen causes
damage to the host. These include, but are not limited to:

\begin{itemize}
    \item \textbf{Adhesins and pili} — enable the bacterium to attach to host tissues.
    \item \textbf{Toxins (exotoxins and endotoxins)} — disrupt host cell function or
          trigger damaging immune responses.
    \item \textbf{Secretion systems (Type I–VI)} — inject effector proteins directly
          into host cells to subvert immune defences.
    \item \textbf{Iron acquisition systems} — allow the pathogen to scavenge iron from
          the host, which is essential for growth.
    \item \textbf{Capsules and immune evasion factors} — protect the bacterium from
          phagocytosis and complement-mediated killing.
    \item \textbf{Biofilm formation} — enables persistent colonisation on tissues and
          implanted medical devices.
\end{itemize}

Because virulence is encoded across many, often distant genomic loci, and because
individual genes may be horizontally transferred between species on mobile genetic
elements, detecting pathogenicity requires a genome-wide analysis rather than
the identification of a single marker gene~\cite{TODO_HGT_virulence}.

\section{Virulence Factor and Antibiotic Resistance Databases}

Two curated public databases underpin the computational detection of virulence
and resistance genes:

\textbf{VFDB (Virulence Factor Database)}~\cite{TODO_VFDB_citation} is a manually curated
repository of virulence-associated sequences from experimentally verified pathogenic
organisms. It is organised into functional categories including adherence, motility,
toxin production, secretion systems, iron acquisition, immune evasion, biofilm, and
invasion. In this work, VFDB was used as the primary reference for virulence gene
detection.
% TODO: Add exact download date and sequence count here, e.g.:
% The version used in this work was downloaded on [DATE] and contained
% [N] protein sequences across [M] functional categories.

\textbf{CARD (Comprehensive Antibiotic Resistance Database)}~\cite{TODO_CARD_citation}
provides a curated collection of antibiotic resistance genes (ARGs) and their
associated mechanisms. CARD is used in this work to assess the antimicrobial
resistance potential of the sequenced microbiome, which is a key component of
clinical pathogenicity assessment.

\section{Computational Detection of Pathogenicity from Nanopore Data}

The bioinformatics workflow for detecting virulence genes from Nanopore sequencing
data involves several steps. First, gene-coding regions are predicted on the raw
long reads using tools such as Prodigal~\cite{TODO_prodigal} or its Python
bindings, Pyrodigal. Predicted coding sequences are then translated into protein
sequences and aligned against the VFDB protein database using DIAMOND
\textit{blastx}~\cite{TODO_diamond}, which natively handles the translation step
from nucleotide queries to protein references.

The resulting alignments are evaluated according to several quality metrics:
percent identity, alignment length, query coverage, and bit score. Thresholds on
these metrics are selected based on the observed score distributions (see
Chapter~\ref{ch:methods}) to distinguish biologically meaningful hits from
spurious low-quality alignments.

Virulence hits are then grouped into functional categories and quantified using
TPM (Transcripts Per Million) normalisation, which allows abundance comparisons
across samples with different total sequencing depths~\cite{TODO_TPM}.

\section{Aims of This Work}

The primary aim of this thesis is to develop a computational pipeline that:

\begin{enumerate}
    \item Accepts Nanopore long-read sequencing data (FASTA format) as input.
    \item Detects and quantifies virulence-associated genes by alignment against VFDB.
    \item Detects and quantifies antibiotic resistance genes by alignment against CARD.
    \item Classifies the detected genes into functional pathogenicity categories.
    \item Assigns a clinically interpretable risk level to each sample.
    \item Enables comparative analysis across multiple samples to assess changes in
          virulence content over time or across conditions.
\end{enumerate}

A secondary aim is to explore the feasibility of using large language models (LLMs)
and protein language models as an alternative, alignment-free approach to virulence
classification, and to compare its performance against the alignment-based pipeline.

The remainder of this thesis is structured as follows.
Chapter~\ref{ch:methods} describes the materials and methods.
% TODO: add remaining chapter descriptions once you have the structure in your main.tex
%  - Replace every TODO comment with the real content
%  - Add your bibliography entries (BibTeX keys are shown next to \cite{})
% ============================================================

\chapter{Introduction}

\section{Clinical Microbiology and Infectious Disease Diagnosis}

Infectious diseases remain a leading cause of morbidity and mortality worldwide,
accounting for a significant proportion of hospital admissions and intensive care
unit stays~\cite{TODO_WHO_infectious_disease}. Rapid and accurate identification
of the causative pathogen is essential for guiding antibiotic therapy, preventing
the spread of infection, and improving patient outcomes.

Traditional diagnostic workflows in clinical microbiology rely on culture-based
methods, which, despite their established reliability, suffer from several
limitations. Bacterial culture typically requires 24 to 72 hours to yield a
result, and many clinically relevant organisms are fastidious or entirely
unculturable under standard laboratory conditions~\cite{TODO_culture_limitations}.
During this waiting period, clinicians are often forced to prescribe broad-spectrum
antibiotics empirically, a practice that contributes to the global problem of
antimicrobial resistance (AMR)~\cite{TODO_AMR_global}.

\section{Next-Generation Sequencing in Clinical Practice}

Next-generation sequencing (NGS) technologies have transformed microbiology by
enabling the direct sequencing of genetic material from clinical samples, bypassing
the need for culture. Metagenomic sequencing in particular allows the simultaneous
detection of bacteria, viruses, fungi, and parasites from a single sequencing run,
providing a comprehensive view of the microbial content of a sample~\cite{TODO_metagenomics_review}.

Short-read sequencing platforms, such as Illumina, dominate the current NGS
landscape. However, the short read lengths produced by these platforms
(typically 150--300 base pairs) pose challenges for resolving complex genomic
regions, mobile genetic elements, and structural variants that are often associated
with virulence and antibiotic resistance~\cite{TODO_shortread_limitations}.

\section{Oxford Nanopore Sequencing}

Oxford Nanopore Technology (ONT) is a third-generation sequencing platform that
produces long reads — typically ranging from several kilobases to over one megabase
in length — by measuring changes in ionic current as a DNA or RNA strand passes
through a nanoscale protein pore~\cite{TODO_ONT_original}. This approach confers
several advantages that are particularly relevant in a clinical context:

\begin{itemize}
    \item \textbf{Speed.} Sequencing results can be obtained within hours of sample
          collection, enabling real-time clinical decision-making~\cite{TODO_ONT_speed}.
    \item \textbf{Long reads.} Read lengths of tens of kilobases allow complete
          assembly of plasmids and mobile genetic elements carrying resistance genes,
          which are often fragmented or misassembled by short-read technologies.
    \item \textbf{Portability and low cost.} The MinION device, weighing less than
          100 grams, can be operated at the point of care, in field settings, or in
          resource-limited environments where large sequencing infrastructure is
          unavailable~\cite{TODO_MinION_pointofcare}.
\end{itemize}

Despite these advantages, Nanopore sequencing is characterised by a higher per-base
error rate (approximately 5--10\% in recent chemistry generations) compared to
Illumina~\cite{TODO_nanopore_accuracy}. This must be accounted for in downstream
bioinformatics analyses, for example by using alignment parameters tuned for
long reads with elevated mismatch rates.

\section{Pathogenicity and Virulence}

Pathogenicity is defined as the ability of a microorganism to cause disease in a
susceptible host. It is not a binary property but rather a complex phenotype
determined by the combined action of multiple genes distributed across the
genome~\cite{TODO_pathogenicity_definition}.

Virulence factors are the molecular mechanisms through which a pathogen causes
damage to the host. These include, but are not limited to:

\begin{itemize}
    \item \textbf{Adhesins and pili} — enable the bacterium to attach to host tissues.
    \item \textbf{Toxins (exotoxins and endotoxins)} — disrupt host cell function or
          trigger damaging immune responses.
    \item \textbf{Secretion systems (Type I–VI)} — inject effector proteins directly
          into host cells to subvert immune defences.
    \item \textbf{Iron acquisition systems} — allow the pathogen to scavenge iron from
          the host, which is essential for growth.
    \item \textbf{Capsules and immune evasion factors} — protect the bacterium from
          phagocytosis and complement-mediated killing.
    \item \textbf{Biofilm formation} — enables persistent colonisation on tissues and
          implanted medical devices.
\end{itemize}

Because virulence is encoded across many, often distant genomic loci, and because
individual genes may be horizontally transferred between species on mobile genetic
elements, detecting pathogenicity requires a genome-wide analysis rather than
the identification of a single marker gene~\cite{TODO_HGT_virulence}.

\section{Virulence Factor and Antibiotic Resistance Databases}

Two curated public databases underpin the computational detection of virulence
and resistance genes:

\textbf{VFDB (Virulence Factor Database)}~\cite{TODO_VFDB_citation} is a manually curated
repository of virulence-associated sequences from experimentally verified pathogenic
organisms. It is organised into functional categories including adherence, motility,
toxin production, secretion systems, iron acquisition, immune evasion, biofilm, and
invasion. In this work, VFDB was used as the primary reference for virulence gene
detection.
% TODO: Add exact download date and sequence count here, e.g.:
% The version used in this work was downloaded on [DATE] and contained
% [N] protein sequences across [M] functional categories.

\textbf{CARD (Comprehensive Antibiotic Resistance Database)}~\cite{TODO_CARD_citation}
provides a curated collection of antibiotic resistance genes (ARGs) and their
associated mechanisms. CARD is used in this work to assess the antimicrobial
resistance potential of the sequenced microbiome, which is a key component of
clinical pathogenicity assessment.

\section{Computational Detection of Pathogenicity from Nanopore Data}

The bioinformatics workflow for detecting virulence genes from Nanopore sequencing
data involves several steps. First, gene-coding regions are predicted on the raw
long reads using tools such as Prodigal~\cite{TODO_prodigal} or its Python
bindings, Pyrodigal. Predicted coding sequences are then translated into protein
sequences and aligned against the VFDB protein database using DIAMOND
\textit{blastx}~\cite{TODO_diamond}, which natively handles the translation step
from nucleotide queries to protein references.

The resulting alignments are evaluated according to several quality metrics:
percent identity, alignment length, query coverage, and bit score. Thresholds on
these metrics are selected based on the observed score distributions (see
Chapter~\ref{ch:methods}) to distinguish biologically meaningful hits from
spurious low-quality alignments.

Virulence hits are then grouped into functional categories and quantified using
TPM (Transcripts Per Million) normalisation, which allows abundance comparisons
across samples with different total sequencing depths~\cite{TODO_TPM}.

\section{Aims of This Work}

The primary aim of this thesis is to develop a computational pipeline that:

\begin{enumerate}
    \item Accepts Nanopore long-read sequencing data (FASTA format) as input.
    \item Detects and quantifies virulence-associated genes by alignment against VFDB.
    \item Detects and quantifies antibiotic resistance genes by alignment against CARD.
    \item Classifies the detected genes into functional pathogenicity categories.
    \item Assigns a clinically interpretable risk level to each sample.
    \item Enables comparative analysis across multiple samples to assess changes in
          virulence content over time or across conditions.
\end{enumerate}

A secondary aim is to explore the feasibility of using large language models (LLMs)
and protein language models as an alternative, alignment-free approach to virulence
classification, and to compare its performance against the alignment-based pipeline.

The remainder of this thesis is structured as follows.
Chapter~\ref{ch:methods} describes the materials and methods.
% TODO: add remaining chapter descriptions once you have the structure in your main.tex
%  - Replace every TODO comment with the real content
%  - Add your bibliography entries (BibTeX keys are shown next to \cite{})
% ============================================================

\chapter{Introduction}

\section{Clinical Microbiology and Infectious Disease Diagnosis}

Infectious diseases remain a leading cause of morbidity and mortality worldwide,
accounting for a significant proportion of hospital admissions and intensive care
unit stays~\cite{TODO_WHO_infectious_disease}. Rapid and accurate identification
of the causative pathogen is essential for guiding antibiotic therapy, preventing
the spread of infection, and improving patient outcomes.

Traditional diagnostic workflows in clinical microbiology rely on culture-based
methods, which, despite their established reliability, suffer from several
limitations. Bacterial culture typically requires 24 to 72 hours to yield a
result, and many clinically relevant organisms are fastidious or entirely
unculturable under standard laboratory conditions~\cite{TODO_culture_limitations}.
During this waiting period, clinicians are often forced to prescribe broad-spectrum
antibiotics empirically, a practice that contributes to the global problem of
antimicrobial resistance (AMR)~\cite{TODO_AMR_global}.

\section{Next-Generation Sequencing in Clinical Practice}

Next-generation sequencing (NGS) technologies have transformed microbiology by
enabling the direct sequencing of genetic material from clinical samples, bypassing
the need for culture. Metagenomic sequencing in particular allows the simultaneous
detection of bacteria, viruses, fungi, and parasites from a single sequencing run,
providing a comprehensive view of the microbial content of a sample~\cite{TODO_metagenomics_review}.

Short-read sequencing platforms, such as Illumina, dominate the current NGS
landscape. However, the short read lengths produced by these platforms
(typically 150--300 base pairs) pose challenges for resolving complex genomic
regions, mobile genetic elements, and structural variants that are often associated
with virulence and antibiotic resistance~\cite{TODO_shortread_limitations}.

\section{Oxford Nanopore Sequencing}

Oxford Nanopore Technology (ONT) is a third-generation sequencing platform that
produces long reads — typically ranging from several kilobases to over one megabase
in length — by measuring changes in ionic current as a DNA or RNA strand passes
through a nanoscale protein pore~\cite{TODO_ONT_original}. This approach confers
several advantages that are particularly relevant in a clinical context:

\begin{itemize}
    \item \textbf{Speed.} Sequencing results can be obtained within hours of sample
          collection, enabling real-time clinical decision-making~\cite{TODO_ONT_speed}.
    \item \textbf{Long reads.} Read lengths of tens of kilobases allow complete
          assembly of plasmids and mobile genetic elements carrying resistance genes,
          which are often fragmented or misassembled by short-read technologies.
    \item \textbf{Portability and low cost.} The MinION device, weighing less than
          100 grams, can be operated at the point of care, in field settings, or in
          resource-limited environments where large sequencing infrastructure is
          unavailable~\cite{TODO_MinION_pointofcare}.
\end{itemize}

Despite these advantages, Nanopore sequencing is characterised by a higher per-base
error rate (approximately 5--10\% in recent chemistry generations) compared to
Illumina~\cite{TODO_nanopore_accuracy}. This must be accounted for in downstream
bioinformatics analyses, for example by using alignment parameters tuned for
long reads with elevated mismatch rates.

\section{Pathogenicity and Virulence}

Pathogenicity is defined as the ability of a microorganism to cause disease in a
susceptible host. It is not a binary property but rather a complex phenotype
determined by the combined action of multiple genes distributed across the
genome~\cite{TODO_pathogenicity_definition}.

Virulence factors are the molecular mechanisms through which a pathogen causes
damage to the host. These include, but are not limited to:

\begin{itemize}
    \item \textbf{Adhesins and pili} — enable the bacterium to attach to host tissues.
    \item \textbf{Toxins (exotoxins and endotoxins)} — disrupt host cell function or
          trigger damaging immune responses.
    \item \textbf{Secretion systems (Type I–VI)} — inject effector proteins directly
          into host cells to subvert immune defences.
    \item \textbf{Iron acquisition systems} — allow the pathogen to scavenge iron from
          the host, which is essential for growth.
    \item \textbf{Capsules and immune evasion factors} — protect the bacterium from
          phagocytosis and complement-mediated killing.
    \item \textbf{Biofilm formation} — enables persistent colonisation on tissues and
          implanted medical devices.
\end{itemize}

Because virulence is encoded across many, often distant genomic loci, and because
individual genes may be horizontally transferred between species on mobile genetic
elements, detecting pathogenicity requires a genome-wide analysis rather than
the identification of a single marker gene~\cite{TODO_HGT_virulence}.

\section{Virulence Factor and Antibiotic Resistance Databases}

Two curated public databases underpin the computational detection of virulence
and resistance genes:

\textbf{VFDB (Virulence Factor Database)}~\cite{TODO_VFDB_citation} is a manually curated
repository of virulence-associated sequences from experimentally verified pathogenic
organisms. It is organised into functional categories including adherence, motility,
toxin production, secretion systems, iron acquisition, immune evasion, biofilm, and
invasion. In this work, VFDB was used as the primary reference for virulence gene
detection.
% TODO: Add exact download date and sequence count here, e.g.:
% The version used in this work was downloaded on [DATE] and contained
% [N] protein sequences across [M] functional categories.

\textbf{CARD (Comprehensive Antibiotic Resistance Database)}~\cite{TODO_CARD_citation}
provides a curated collection of antibiotic resistance genes (ARGs) and their
associated mechanisms. CARD is used in this work to assess the antimicrobial
resistance potential of the sequenced microbiome, which is a key component of
clinical pathogenicity assessment.

\section{Computational Detection of Pathogenicity from Nanopore Data}

The bioinformatics workflow for detecting virulence genes from Nanopore sequencing
data involves several steps. First, gene-coding regions are predicted on the raw
long reads using tools such as Prodigal~\cite{TODO_prodigal} or its Python
bindings, Pyrodigal. Predicted coding sequences are then translated into protein
sequences and aligned against the VFDB protein database using DIAMOND
\textit{blastx}~\cite{TODO_diamond}, which natively handles the translation step
from nucleotide queries to protein references.

The resulting alignments are evaluated according to several quality metrics:
percent identity, alignment length, query coverage, and bit score. Thresholds on
these metrics are selected based on the observed score distributions (see
Chapter~\ref{ch:methods}) to distinguish biologically meaningful hits from
spurious low-quality alignments.

Virulence hits are then grouped into functional categories and quantified using
TPM (Transcripts Per Million) normalisation, which allows abundance comparisons
across samples with different total sequencing depths~\cite{TODO_TPM}.

\section{Aims of This Work}

The primary aim of this thesis is to develop a computational pipeline that:

\begin{enumerate}
    \item Accepts Nanopore long-read sequencing data (FASTA format) as input.
    \item Detects and quantifies virulence-associated genes by alignment against VFDB.
    \item Detects and quantifies antibiotic resistance genes by alignment against CARD.
    \item Classifies the detected genes into functional pathogenicity categories.
    \item Assigns a clinically interpretable risk level to each sample.
    \item Enables comparative analysis across multiple samples to assess changes in
          virulence content over time or across conditions.
\end{enumerate}

A secondary aim is to explore the feasibility of using large language models (LLMs)
and protein language models as an alternative, alignment-free approach to virulence
classification, and to compare its performance against the alignment-based pipeline.

The remainder of this thesis is structured as follows.
Chapter~\ref{ch:methods} describes the materials and methods.
% TODO: add remaining chapter descriptions once you have the structure